\scanpage{21}
\begin{center}4 -- 24 --\end{center}
\hfill n° 392

D'autre part, $(A:\underline a)_{\underline p}=A_{\underline p}:\underline a A_{\underline p}=\underline p^{-n_{\underline p}}A_{\underline p}$, donc $\ell_{\underline p}((A:\underline a)/A)=n_{\underline p}$, et on conclut de la même façon.

6. \uline{Invariant relatif de deux réseaux.}

Les notations et hypothèses sont celles des n\textsuperscript{os} 2 à 5. Soient $V$ un espace vectoriel de dimension finie $n$ sur $K$, $M$ un réseau de $V$ par rapport à $A$. Soit $W$ la puissance extérieure $\overset{n}{\bigwedge}V$, qui est un espace vectoriel de \uline{dimension} 1 sur $K$, et désignons par $M_W$ le réseau de $W$ engendré par l'image de $M^n$ par l'application canonique $V^n\to\overset{n}{\bigwedge}V$, (n° 1, prop. 3, (iii); on notera que $M_W$ n'est pas nécessairement isomorphe à $\overset{n}{\bigwedge}M$, (\uline{Alg.}, chap. III, § 5, exerc. 9)). Si $e$ est une base de $W$ sur $K$, on peut donc écrire $M_W=\underline a.e$, où $\underline a$ est un idéal fractionnaire $\ne0$ de $A$.

Soit $M'$ un second réseau de $V$, et posons $M'_W=\underline a'.e$, où $\underline a'$ est un idéal fractionnaire $\ne0$ de $A$; le diviseur $\operatorname{div}(\underline a)-\operatorname{div}(\underline a')$ \uline{ne dépend pas du choix de la base} $e$ \uline{de} $W$, $\underline a$ et $\underline a'$ étant multipliés par un même élément de $K^*$ quand on change de base; nous poserons $\chi(M,M')=\operatorname{div}(\underline a)-\operatorname{div}(\underline a')$ et nous dirons que ce diviseur est l'\uline{invariant relatif de} $M'$ \uline{par rapport à} $M$. Il est clair que si $M,M',M''$ sont trois réseaux de $V$, on a:
\[
(3)\qquad\chi(M,M')+\chi(M',M'')+\chi(M'',M)=0
\]
\[
(4)\qquad\chi(M,M')+\chi(M',M)=0.
\]

Pour tout $\underline p\in P$, il résulte aussitôt des définitions que l'on a $(M_W)_{\underline p}=(M_{\underline p})_W$; en outre, $M_{\underline p}$ étant alors un $A_{\underline p}$-module \uline{libre} puisque $A_{\underline p}$ est principal, une base de $M_{\underline p}$ sur $A_{\underline p}$ est une base de $V$ sur $K$, donc $(M_{\underline p})_W=\overset{n}{\bigwedge}(M_{\underline p})$, (chap. II, § 2, n° 8), et l'idéal fractionnaire $\underline a_{\underline p}=\underline a A_{\underline p}$ est principal. Si on pose $\underline a_{\underline p}=\underline p^{n_{\underline p}}A_{\underline p}$, $\underline a'_{\underline p}=\underline p^{n'_{\underline p}}A_{\underline p}$, on a donc:
\[
\chi(M,M')=\sum_{\underline p\in P}(n_{\underline p}-n'_{\underline p}).\underline p\ ,
\]
\mspage{20}
